\section*{Относительная ориентация верхнего и нижнего секторов}

После факторизации по остаточной группе остаются два относительных класса:
совпадающие и ортогональные представители осевой орбиты. Для минимального
оператора
\[
 M_s=L(w_3,w_1)+gA_{n_s}
\]
совпадающий класс даёт единичное смешивание. Ортогональный класс даёт только
вещественное вращение первых двух характеров,
\[
 \tan\theta_{\rm mix}=\frac{g(1-c_*)}{2(w_3-w_1)},
\]
тогда как третье поколение не смешивается. Коэффициент (g) не выведен,
поэтому полный беспараметрический вывод CKM не достигнут.